\section{Related work}
\label{sec:related}

\paragraph{Generalised card representations.}
Bertram, F\"urnkranz, and M\"uller~\citep{bertram2024cards} established the
zero-shot drafting problem, and their paper remains the only published
quantitative result on it. They replace one-hot card identities with
concatenated feature vectors: hand-engineered attributes plus a
Sentence-BERT embedding~\citep{reimers2019sbert} of card text (1306-d), a
learned autoencoding of card images (1024-d), and 16 post-release usage
statistics from \seventeenlands{}. This lives inside a
contextual-preference-ranking model trained with a triplet loss. Two of
their findings shape our design: feature
representations match one-hot accuracy within a set (nothing is lost by
generalising), and multi-set pretraining, not the representation alone, is
what buys transfer. Pretrained on 13 sets and evaluated on held-out BRO,
their best representation reaches 55.44\% pick agreement. That
representation includes the usage-statistics block, which is computed from
human play after a set's release; the number is therefore an anchor for
zero-shot \emph{card generalisation} rather than for day-one deployment,
where such statistics cannot exist. The same paper's day-one-feasible
representation (features only) reaches 33.6--35.6\% on unseen sets under
single-set training. No multi-set features-only number is published.
Filling that gap (multi-set training, public card information only) is
this paper's headline experiment, and we adopt their held-out BRO as a
development set so the comparison is on the same holdout. A companion
paper~\citep{bertram2024infonce} improves the within-set triplet objective
with an adapted InfoNCE loss but reports no unseen-set results.

\paragraph{Within-set pick models.}
Ward et al.~\citep{ward2021drafting} compare heuristic and neural drafters
on M19 and report the two most useful reference points below the
data-driven regime: an expert-tuned rating bot at 44.5\% top-1 agreement
and a random agent at 22\%. Their neural drafter (48.7\%) and its
successors, the statistical-drafting MLPs deployed at
statisticaldrafting.com~\citep{brooks2024statistical} (${\sim}70\%$) and
Czerner's puder transformer~\citep{czerner2025puder} (71\%), define the
within-set state of the art, and the statistical-drafting architecture is
the direct lineage of the per-set ceiling models we compare against.
All are one-hot models over a fixed card vocabulary; none can score an
unseen set. The deployment gap is explicit: statistical-drafting documents
that no model can ship until \seventeenlands{} data appears, roughly two
weeks post-release~\citep{brooks2024statistical}, and Draftsim's day-one
ratings are written by human experts~\citep{draftsimratings}.

\paragraph{LLMs as drafters.}
UrzaGPT~\citep{bertram2025urzagpt} evaluates large language models on draft
picks. Zero-shot GPT-4o reaches 43\%, narrowly above the hour-zero rarity
heuristic we measure (\S\ref{sec:results}). We still mark it non-day-one
in Table~\ref{tab:baselines}: NEO shipped years before GPT-4o's training
cutoff, so its pretraining corpus cannot be shown free of the set's own
post-release strategy discussion, a leak a frozen feature-only encoder
cannot have. Providing full rules text in the prompt lowers GPT-4o's
accuracy further. Open 7--8B models fail to pick legal cards until
LoRA-tuned on a million picks from the target set, after which they still
trail small supervised MLPs at orders of magnitude more inference cost. We
take the implied guidance: use a frozen text \emph{embedding} as one
feature block, not a language model as the policy.

\paragraph{Text-grounded transfer in card games.}
Cardsformer~\citep{xia2023cardsformer} encodes Hearthstone card
descriptions with a language model so a gameplay policy generalises to
unseen cards. It is the closest cross-game precedent for text-mediated
transfer, applied to gameplay rather than drafting.
